%==============================================================================%
% Chapter 1 – Introduction
%==============================================================================%

\chapter{Introduction}
\label{ch:intro}

\section{Motivation}

The threat posed by large-scale quantum computers to today's public-key
infrastructure is no longer a theoretical curiosity.  Shor's algorithm reduces
the security of RSA, DSA, and elliptic-curve schemes to polynomial time once a
sufficiently powerful quantum computer is available~\cite{Shor1994}.  In
response, the U.S.\ National Institute of Standards and Technology (NIST) ran
a multi-year post-quantum cryptography standardisation project, selecting four
algorithms for standardisation in 2022--2024~\cite{NIST2022}, while the Korean
Cryptography Module Validation Programme (KPQC) is pursuing its own national
standards~\cite{KPQC2022}.

\haetae{}~\cite{HAETAE2022} is one of the KPQC signature finalists.  It is a
lattice-based scheme whose security reduces to the Decisional Module
Short Integer Solution (DMSIS) hardness assumption over module lattices.  Like
all modern lattice-based schemes, \haetae{} relies on efficient polynomial
multiplication modulo $X^{256}+1$ in the ring $\Z_q[X]/(X^{256}+1)$, which
in practice is implemented via the \emph{Number-Theoretic Transform} (NTT).
The NTT reduces an $O(n^2)$ schoolbook convolution to $O(n \log n)$ operations
and is the dominant computational primitive in the scheme.

For a cryptographic standard to be trustworthy, correctness proofs cannot stop
at pen-and-paper mathematics.  History has shown that hand-verified
implementations harbour subtle bugs: off-by-one errors in loop bounds,
incorrect twiddle factor tables, wrong modular arithmetic, or mishandled
Montgomery scaling can silently produce wrong outputs while passing standard
test vectors~\cite{BerBuc2017}.  This dissertation takes the position that
\emph{machine-checked} formal proofs, linking every line of executable code to
its abstract mathematical specification, are the appropriate gold standard for
cryptographic primitives in post-quantum standards.

\section{Problem Statement}

We address the following problem:

\begin{quote}
  \textit{Prove, in a machine-checkable formal system, that the \jasmin{}
  implementation of the NTT and inverse-NTT in \haetae{} is functionally
  correct with respect to the abstract mathematical specification of the
  negacyclic NTT over $\Z_q[X]/(X^{256}+1)$.}
\end{quote}

The challenge is to bridge three distinct abstraction levels:

\begin{enumerate}
  \item \textbf{Low-level word arithmetic.}  The \jasmin{} source is compiled
        to x86-64 assembly and works with 32-bit and 64-bit machine words.
        Correctness at this level means proving arithmetic identities modulo
        $2^{32}$ and establishing bounds (e.g., that an intermediate product
        fits in a 64-bit register without overflow).

  \item \textbf{Algorithmic (loop) level.}  The Cooley--Tukey butterfly
        algorithm is expressed as a nested three-loop structure with eight
        stages.  Correctness here means showing that after all loops terminate,
        the array holds the intended Fourier coefficients, using loop invariants
        expressed in terms of partial NTT computations.

  \item \textbf{Mathematical (spectral) level.}  The NTT is a linear map
        $\Z_q^{256} \to \Z_q^{256}$ defined by evaluation at specific powers
        of a primitive $512$th root of unity.  Correctness here means proving
        that the algorithmic output equals the evaluation formula, and that the
        inverse NTT is its exact algebraic inverse.
\end{enumerate}

Each layer requires different mathematical machinery: modular arithmetic
and bit-vector identities at level~1, loop invariants and array
relational assertions at level~2, and field theory with primitive roots at
level~3.  Composing all three in a single proof assistant is the central
technical challenge.

\section{Tools and Methodology}

We use two mature high-assurance tools:

\paragraph{\jasmin{}~\cite{Almeida2017}.}
A domain-specific programming language and compiler designed for
high-assurance cryptographic implementations.  \jasmin{} programs are compiled
to provably correct x86-64 assembly and simultaneously admit extraction into
the \easycrypt{} proof assistant as first-class program modules.  The language
provides explicit register and memory annotations that remove ambiguity from
the compiler's optimisation decisions.

\paragraph{\easycrypt{}~\cite{Barthe2011,Barthe2013}.}
A proof assistant built around the probabilistic Relational Hoare Logic
($\pRHL$) calculus, designed specifically for reasoning about cryptographic
constructions.  An \easycrypt{} \emph{program logic} allows Hoare-style
reasoning over both deterministic and probabilistic procedures.  The ambient
logic is a dependent type theory supporting classical higher-order reasoning,
enabling rich algebraic developments (ring theory, field theory, group theory)
within the same tool.

Our proof strategy follows the \emph{refinement} paradigm:
\begin{align*}
  \text{Jasmin} &\xrightarrow{\text{extraction}}
  \text{EC module (RefJasminNTT)} \xrightarrow{\text{pRHL equivalence}}
  \text{EC imperative spec (NTT\_Fq)} \\
  &\xrightarrow{\text{loop invariant}}
  \text{EC abstract spec (NTTFullSpec/Algebra)}
\end{align*}
Each arrow represents a machine-checked proof step; the composition yields
the end-to-end theorem stated in \ec{NTTEndToEnd.ec}.

\section{Main Contributions}

\begin{description}
  \item[\textbf{C1 -- Montgomery correctness.}]
    A machine-checked proof that \ec{montgomery\_reduce} computes
    $a \cdot R^{-1} \bmod q$ for all inputs in the relevant working range,
    with $R = 2^{32}$, $q = 64513$, explicitly verified bit-width bounds,
    and signed-integer arithmetic lemmas absent from the \easycrypt{} standard
    library (\Cref{ch:montgomery}).

  \item[\textbf{C2 -- Imperative NTT specification.}]
    A fully proved \easycrypt{} model of the forward and inverse NTT loops in
    \ec{NTT\_Fq.ec}, including twiddle-factor correctness, coefficient-bound
    invariants (16-bit input $\to$ 24-bit output for forward NTT; 16-bit
    preservation for inverse NTT), and termination witnesses
    (\Cref{ch:ntt_spec}).

  \item[\textbf{C3 -- Abstract spectral specification.}]
    An algebraic library \ec{NTTFullAlgebra.ec} establishing properties of the
    primitive root $\zeta_0 = 426$ in $\Zq$, bit-reversal permutation lemmas,
    and a butterfly-stage induction connecting the iterative algorithm to the
    closed-form evaluation formula (\Cref{ch:algebra}).

  \item[\textbf{C4 -- Jasmin refinement.}]
    A pRHL proof in \ec{RefJasminNTT.ec} showing that the extracted \jasmin{}
    procedures are observationally equivalent to the imperative specifications
    up to array representation (\Cref{ch:refinement}).

  \item[\textbf{C5 -- End-to-end correctness.}]
    Top-level theorems \ec{haetae\_jasmin\_ntt\_correct} and
    \ec{haetae\_jasmin\_invntt\_correct} in \ec{NTTEndToEnd.ec}, together
    with the round-trip theorem $\INTT \circ \NTT = \op{id}$
    (\Cref{ch:end_to_end}).
\end{description}

\section{Structure of the Dissertation}

The remainder of this dissertation is organised as follows.

\begin{description}
  \item[\Cref{ch:background}] reviews the relevant background: lattice
        cryptography, the NTT, post-quantum signature schemes, and the
        theoretical foundations of formal verification with \easycrypt{} and
        \jasmin{}.

  \item[\Cref{ch:haetae_ntt}] describes the \haetae{} scheme, its NTT
        parameters, and the reference C implementation that serves as the
        algorithmic ground truth.

  \item[\Cref{ch:formal_framework}] presents the overall formal framework:
        the three-layer abstraction stack and the proof methodology.

  \item[\Cref{ch:field_ring}] covers the formalization of the field
        $\Zq$ and the polynomial ring $\Rq$ in \easycrypt{}.

  \item[\Cref{ch:montgomery}] details the Montgomery reduction
        formalization and its correctness proof.

  \item[\Cref{ch:ntt_spec}] presents the imperative NTT specification in
        \ec{NTT\_Fq.ec}, its invariants, and its termination proof.

  \item[\Cref{ch:algebra}] develops the algebraic infrastructure
        in \ec{NTTFullAlgebra.ec}: primitive root properties, bit-reversal
        lemmas, and stage decompositions.

  \item[\Cref{ch:jasmin_impl}] describes the \jasmin{} implementation in
        detail.

  \item[\Cref{ch:refinement}] presents the Jasmin-to-EasyCrypt refinement
        proof in \ec{RefJasminNTT.ec}.

  \item[\Cref{ch:end_to_end}] assembles all pieces into the end-to-end
        correctness theorems.

  \item[\Cref{ch:conclusion}] discusses limitations, remaining obligations,
        related and future work.
\end{description}
